\documentclass[12pt,a4paper]{article}

\usepackage[T2A]{fontenc}
\usepackage[utf8]{inputenc}
\usepackage[russian]{babel}
\usepackage{amsmath,amssymb,amsthm}
\usepackage{array}
\usepackage{booktabs}
\usepackage{longtable}
\usepackage{enumitem}
\usepackage{geometry}
\usepackage{hyperref}
\usepackage{xcolor}

\geometry{left=25mm,right=15mm,top=20mm,bottom=20mm}

\hypersetup{
    colorlinks=true,
    linkcolor=blue,
    urlcolor=blue
}

\title{Билет №63 \\ \small UML CASE-системы: описание функциональности, примеры известных систем. (СпБГУ)}
\author{Сериков Владислав Александрович}
\date{16 мая 2026}

\begin{document}

\maketitle
\tableofcontents
\newpage

\section{Базовые понятия}

\subsection{CASE-системы}

\textbf{CASE} --- это сокращение от \textit{Computer-Aided Software Engineering}, то есть ``автоматизированная разработка программного обеспечения''.

\textbf{CASE-система} --- программная система, предназначенная для поддержки процессов анализа, проектирования, разработки, тестирования, документирования, сопровождения и управления программными системами.

CASE-системы применяются для:
\begin{itemize}
    \item формализации требований;
    \item построения моделей предметной области и программной архитектуры;
    \item автоматической генерации документации;
    \item генерации исходного кода по моделям;
    \item обратного проектирования, то есть восстановления моделей по существующему коду;
    \item обеспечения согласованности моделей и артефактов разработки;
    \item коллективной работы над проектом;
    \item трассировки требований, проектных решений, кода и тестов.
\end{itemize}

\subsection{UML}

\textbf{UML} --- \textit{Unified Modeling Language}, унифицированный язык моделирования.

\textbf{UML} --- это стандартизованный графический язык, предназначенный для спецификации, визуализации, проектирования и документирования систем, прежде всего программных и программно-аппаратных.

UML не является языком программирования. Он описывает систему на более высоком уровне абстракции: через объекты, классы, компоненты, состояния, взаимодействия, сценарии использования и другие сущности.

\subsection{UML CASE-система}

\textbf{UML CASE-система} --- CASE-система, основным средством моделирования в которой является язык UML.

Такая система предоставляет пользователю инструменты для создания UML-диаграмм, хранения модели, проверки её корректности, генерации документации, а в более развитых вариантах --- генерации кода, синхронизации модели и кода, командной работы и интеграции с другими средствами разработки.

\section{Место UML CASE-систем в жизненном цикле ПО}

UML CASE-системы могут использоваться практически на всех стадиях жизненного цикла программного обеспечения.

\begin{longtable}{p{0.25\textwidth}p{0.65\textwidth}}
\toprule
\textbf{Стадия} & \textbf{Роль UML CASE-системы} \\
\midrule
Анализ требований &
Построение диаграмм вариантов использования, описание акторов, сценариев, ограничений, бизнес-процессов. \\
\addlinespace

Анализ предметной области &
Построение концептуальных моделей: классы предметной области, связи, обобщения, композиции, ассоциации. \\
\addlinespace

Проектирование архитектуры &
Описание компонентов, подсистем, интерфейсов, слоёв, зависимостей, deployment-структуры. \\
\addlinespace

Детальное проектирование &
Построение диаграмм классов, последовательностей, состояний, активностей, взаимодействий. \\
\addlinespace

Реализация &
Генерация каркаса кода, синхронизация модели и исходного кода, обратное проектирование. \\
\addlinespace

Тестирование &
Получение тестовых сценариев из вариантов использования, трассировка требований к тестам. \\
\addlinespace

Сопровождение &
Документирование существующей системы, анализ зависимостей, восстановление архитектуры по коду. \\
\bottomrule
\end{longtable}

\section{Классификация CASE-систем}

\subsection{По охвату жизненного цикла}

\begin{enumerate}
    \item \textbf{Upper CASE} --- средства, поддерживающие ранние стадии разработки:
    \begin{itemize}
        \item анализ требований;
        \item бизнес-моделирование;
        \item концептуальное проектирование;
        \item архитектурное проектирование.
    \end{itemize}

    \item \textbf{Lower CASE} --- средства, ориентированные на поздние стадии:
    \begin{itemize}
        \item детальное проектирование;
        \item генерация кода;
        \item тестирование;
        \item сопровождение;
        \item обратное проектирование.
    \end{itemize}

    \item \textbf{Integrated CASE} --- интегрированные CASE-системы, поддерживающие полный цикл:
    \begin{itemize}
        \item требования;
        \item проектирование;
        \item кодирование;
        \item тестирование;
        \item документацию;
        \item управление изменениями.
    \end{itemize}
\end{enumerate}

\subsection{По назначению}

\begin{itemize}
    \item \textbf{средства бизнес-моделирования} --- описание бизнес-процессов, организационных структур, потоков данных;
    \item \textbf{средства объектно-ориентированного проектирования} --- построение UML-моделей программных систем;
    \item \textbf{средства системной инженерии} --- моделирование сложных программно-аппаратных систем, часто с поддержкой SysML;
    \item \textbf{средства архитектурного моделирования} --- описание архитектуры предприятия и программных ландшафтов;
    \item \textbf{средства проектирования баз данных} --- ER-моделирование, генерация DDL, обратное проектирование схем БД;
    \item \textbf{средства генерации документации} --- автоматическое получение проектных документов из модели.
\end{itemize}

\subsection{По степени формальности}

\begin{enumerate}
    \item \textbf{Диаграммные редакторы.}

    Такие инструменты позволяют рисовать диаграммы, но модель обычно хранится как набор графических объектов. Семантическая проверка слабая или отсутствует.

    Примеры:
    \begin{itemize}
        \item Microsoft Visio;
        \item draw.io / diagrams.net;
        \item Lucidchart.
    \end{itemize}

    \item \textbf{Модельно-ориентированные UML-системы.}

    В таких системах диаграмма является представлением внутренней модели. Один и тот же элемент модели может отображаться на разных диаграммах.

    Примеры:
    \begin{itemize}
        \item Enterprise Architect;
        \item Visual Paradigm;
        \item IBM Rational Software Architect;
        \item Modelio;
        \item StarUML.
    \end{itemize}

    \item \textbf{Средства model-driven development.}

    Эти системы поддерживают разработку, управляемую моделями:
    \begin{itemize}
        \item генерацию кода;
        \item трансформацию моделей;
        \item профили UML;
        \item проверку ограничений;
        \item синхронизацию модели и реализации.
    \end{itemize}
\end{enumerate}

\section{Основная функциональность UML CASE-систем}

\subsection{Создание UML-диаграмм}

UML CASE-система должна поддерживать основные виды UML-диаграмм.

\subsubsection{Структурные диаграммы}

\begin{itemize}
    \item \textbf{Диаграмма классов} --- описывает классы, интерфейсы, атрибуты, операции и отношения между типами.
    \item \textbf{Диаграмма объектов} --- показывает экземпляры классов и связи между ними в конкретный момент времени.
    \item \textbf{Диаграмма компонентов} --- описывает программные компоненты и их зависимости.
    \item \textbf{Диаграмма развёртывания} --- показывает физические узлы, устройства, серверы и размещение артефактов.
    \item \textbf{Диаграмма пакетов} --- группирует элементы модели и показывает зависимости между пакетами.
    \item \textbf{Диаграмма композитной структуры} --- описывает внутреннее устройство класса или компонента.
\end{itemize}

\subsubsection{Поведенческие диаграммы}

\begin{itemize}
    \item \textbf{Диаграмма вариантов использования} --- описывает функциональные требования с точки зрения внешних акторов.
    \item \textbf{Диаграмма активностей} --- моделирует поток управления или бизнес-процесс.
    \item \textbf{Диаграмма состояний} --- описывает состояния объекта и переходы между ними.
\end{itemize}

\subsubsection{Диаграммы взаимодействия}

\begin{itemize}
    \item \textbf{Диаграмма последовательностей} --- показывает обмен сообщениями между объектами во времени.
    \item \textbf{Диаграмма коммуникации} --- показывает взаимодействие объектов с акцентом на связях.
    \item \textbf{Диаграмма обзора взаимодействия} --- объединяет элементы диаграмм активностей и взаимодействий.
    \item \textbf{Диаграмма синхронизации} --- описывает изменение состояний объектов во времени.
\end{itemize}

\subsection{Хранение единой модели}

В развитых UML CASE-системах диаграмма не является просто рисунком. Она представляет некоторую часть общей модели.

Например, если класс \texttt{User} присутствует на диаграмме классов и используется на диаграмме последовательностей, то в модели должен существовать один элемент \texttt{User}, имеющий несколько графических представлений.

Преимущества единого репозитория модели:
\begin{itemize}
    \item устранение дублирования;
    \item возможность проверки согласованности;
    \item трассировка связей между элементами;
    \item автоматическая генерация документации;
    \item поддержка командной разработки;
    \item возможность трансформации модели в код.
\end{itemize}

\subsection{Проверка корректности модели}

UML CASE-система может выполнять различные проверки:
\begin{itemize}
    \item наличие неименованных элементов;
    \item корректность связей между элементами;
    \item отсутствие циклических зависимостей в запрещённых местах;
    \item соответствие диаграмм правилам UML;
    \item соответствие модели корпоративным стандартам;
    \item выполнение ограничений, заданных пользователем.
\end{itemize}

Например, отношение наследования должно соединять совместимые классификаторы, а сообщение на диаграмме последовательностей должно соответствовать операции принимающего объекта.

\subsection{Трассировка требований}

\textbf{Трассировка} --- установление и поддержание связей между различными артефактами проекта.

Типовые связи:
\[
\text{требование} \rightarrow \text{вариант использования} \rightarrow \text{класс} \rightarrow \text{метод} \rightarrow \text{тест}.
\]

Трассировка позволяет ответить на вопросы:
\begin{itemize}
    \item какие классы реализуют данное требование;
    \item какие требования покрыты тестами;
    \item какие части системы затронет изменение требования;
    \item какие требования ещё не реализованы;
    \item какие элементы модели не имеют связи с требованиями.
\end{itemize}

\subsection{Генерация документации}

UML CASE-системы позволяют автоматически формировать проектную документацию:
\begin{itemize}
    \item описание требований;
    \item спецификацию вариантов использования;
    \item архитектурное описание;
    \item описание классов и интерфейсов;
    \item описание компонентов;
    \item отчёты о зависимостях;
    \item HTML-, PDF-, DOCX- или RTF-документы.
\end{itemize}

Автоматическая генерация документации важна потому, что документация, написанная вручную, быстро устаревает. Если документация строится из модели, то достаточно поддерживать актуальность модели.

\subsection{Генерация кода}

\textbf{Прямое проектирование} --- получение исходного кода из модели.

Обычно генерация выполняется по диаграммам классов и компонентным моделям.

Из класса UML можно получить:
\begin{itemize}
    \item объявление класса в языке программирования;
    \item поля по атрибутам;
    \item методы по операциям;
    \item интерфейсы;
    \item связи наследования;
    \item ассоциации в виде ссылочных полей;
    \item комментарии из описаний модели.
\end{itemize}

\subsubsection{Минимальный пример генерации кода}

Пусть в UML задан класс:

\[
\texttt{User}
\]

с атрибутами:
\[
\texttt{id: int}, \qquad \texttt{name: String}
\]

и операцией:
\[
\texttt{getName(): String}.
\]

Тогда для Java может быть сгенерирован следующий каркас:

\begin{verbatim}
public class User {
    private int id;
    private String name;

    public String getName() {
        return name;
    }
}
\end{verbatim}

Важно, что CASE-система обычно генерирует не полный прикладной код, а структуру:
\begin{itemize}
    \item классы;
    \item интерфейсы;
    \item сигнатуры методов;
    \item связи между классами;
    \item иногда шаблонные реализации.
\end{itemize}

\subsection{Обратное проектирование}

\textbf{Обратное проектирование} --- построение модели по уже существующим артефактам, чаще всего по исходному коду или базе данных.

Например, CASE-система может прочитать Java-код и построить по нему диаграмму классов.

\subsubsection{Пример}

Исходный код:

\begin{verbatim}
public class Order {
    private Customer customer;
    private List<OrderItem> items;

    public double total() {
        // ...
    }
}
\end{verbatim}

После обратного проектирования можно получить UML-класс \texttt{Order}:
\begin{itemize}
    \item атрибут \texttt{customer: Customer};
    \item атрибут \texttt{items: List<OrderItem>};
    \item операция \texttt{total(): double};
    \item ассоциации с классами \texttt{Customer} и \texttt{OrderItem}.
\end{itemize}

\subsection{Round-trip engineering}

\textbf{Round-trip engineering} --- двусторонняя синхронизация модели и кода.

Идея:
\[
\text{модель} \leftrightarrow \text{код}.
\]

Если разработчик изменяет диаграмму классов, система обновляет код. Если разработчик изменяет код, система обновляет модель.

Это сложная задача, потому что модель обычно абстрактнее кода. Не всякая конструкция языка программирования имеет прямое UML-представление, и не всякая UML-конструкция однозначно отображается в код.

\subsection{Поддержка командной разработки}

Развитые CASE-системы поддерживают:
\begin{itemize}
    \item общий репозиторий моделей;
    \item разграничение прав доступа;
    \item блокировку элементов при редактировании;
    \item версионирование;
    \item сравнение моделей;
    \item слияние изменений;
    \item интеграцию с Git, SVN, TFS и другими системами контроля версий;
    \item публикацию моделей через веб-интерфейс.
\end{itemize}

\subsection{Профили UML и расширяемость}

\textbf{Профиль UML} --- механизм расширения UML для конкретной предметной области.

Профиль обычно включает:
\begin{itemize}
    \item стереотипы;
    \item tagged values, то есть именованные свойства;
    \item ограничения;
    \item специальные иконки и нотации.
\end{itemize}

Примеры профилей и связанных языков:
\begin{itemize}
    \item SysML --- язык системного моделирования, основанный на UML;
    \item MARTE --- профиль для моделирования встроенных и реального времени систем;
    \item SoaML --- профиль для сервисно-ориентированных архитектур.
\end{itemize}

\section{Типовая архитектура UML CASE-системы}

UML CASE-система обычно состоит из следующих подсистем:

\begin{enumerate}
    \item \textbf{Графический редактор диаграмм.}

    Позволяет создавать и редактировать диаграммы, размещать элементы, соединять их отношениями.

    \item \textbf{Репозиторий модели.}

    Хранит элементы модели, связи, свойства, версии, метаданные.

    \item \textbf{Метамодель.}

    Описывает допустимые типы элементов: класс, интерфейс, ассоциация, компонент, актор, состояние и т.д.

    \item \textbf{Механизм проверки модели.}

    Выполняет проверку синтаксических и семантических правил.

    \item \textbf{Генератор кода.}

    Преобразует модель в исходный код на целевом языке.

    \item \textbf{Механизм обратного проектирования.}

    Строит модель по существующему коду или схемам данных.

    \item \textbf{Генератор документации.}

    Формирует отчёты и проектные документы.

    \item \textbf{Средства интеграции.}

    Обеспечивают связь с IDE, системами контроля версий, баг-трекерами, системами управления требованиями.

    \item \textbf{Средства командной работы.}

    Обеспечивают совместное редактирование, роли, права, блокировки и публикацию моделей.
\end{enumerate}

\section{Алгоритмы и типовые процедуры}

\subsection{Алгоритм построения UML-модели в CASE-системе}

\textbf{Цель:} получить согласованную модель программной системы.

\subsubsection*{Шаги}

\begin{enumerate}
    \item \textbf{Определить границы системы.}

    Нужно понять, что входит в систему, а что является внешней средой.

    \item \textbf{Выделить акторов.}

    Актор --- внешняя роль, взаимодействующая с системой.

    \item \textbf{Построить диаграмму вариантов использования.}

    Фиксируются основные функции системы.

    \item \textbf{Описать сценарии вариантов использования.}

    Для каждого варианта использования задаются основной сценарий и альтернативные сценарии.

    \item \textbf{Построить концептуальную диаграмму классов.}

    Выделяются основные сущности предметной области.

    \item \textbf{Построить диаграммы взаимодействия.}

    Для ключевых сценариев строятся диаграммы последовательностей.

    \item \textbf{Уточнить диаграмму классов.}

    Добавляются операции, атрибуты, интерфейсы, зависимости.

    \item \textbf{Описать состояния объектов.}

    Для объектов со сложным жизненным циклом строятся диаграммы состояний.

    \item \textbf{Спроектировать компоненты и развёртывание.}

    Определяется физическая и логическая архитектура.

    \item \textbf{Проверить модель.}

    Выполняется поиск противоречий, неиспользуемых элементов и нарушений правил.

    \item \textbf{Сгенерировать документацию или код.}

    В зависимости от целей проекта модель используется для дальнейших артефактов.
\end{enumerate}

\subsubsection*{Минимальный пример}

Пусть проектируется интернет-магазин.

\begin{enumerate}
    \item Граница системы: веб-приложение интернет-магазина.
    \item Акторы: \texttt{Покупатель}, \texttt{Администратор}, \texttt{Платёжная система}.
    \item Варианты использования:
    \begin{itemize}
        \item \texttt{Просмотреть каталог};
        \item \texttt{Добавить товар в корзину};
        \item \texttt{Оформить заказ};
        \item \texttt{Оплатить заказ};
        \item \texttt{Управлять товарами}.
    \end{itemize}
    \item Классы предметной области:
    \begin{itemize}
        \item \texttt{User};
        \item \texttt{Product};
        \item \texttt{Cart};
        \item \texttt{Order};
        \item \texttt{Payment}.
    \end{itemize}
    \item Для сценария \texttt{Оформить заказ} строится диаграмма последовательностей:
    \[
    \texttt{Customer} \rightarrow \texttt{Cart} \rightarrow \texttt{OrderService} \rightarrow \texttt{PaymentGateway}.
    \]
\end{enumerate}

\subsection{Алгоритм генерации кода по диаграмме классов}

\textbf{Вход:} UML-модель классов.

\textbf{Выход:} каркас исходного кода на выбранном языке.

\subsubsection*{Шаги}

\begin{enumerate}
    \item Для каждого UML-класса создать файл исходного кода.
    \item Для каждого атрибута класса создать поле.
    \item Для каждой операции создать метод с соответствующей сигнатурой.
    \item Для обобщения создать наследование.
    \item Для реализации интерфейса создать конструкцию \texttt{implements} или аналог.
    \item Для ассоциаций создать ссылочные поля или коллекции.
    \item Для кратности \texttt{0..*} или \texttt{1..*} использовать коллекцию.
    \item Для комментариев и описаний модели создать комментарии в коде.
    \item Сохранить полученные файлы.
\end{enumerate}

\subsubsection*{Иллюстрация}

UML-модель:
\[
\texttt{Customer} \; 1 \longrightarrow 0..* \; \texttt{Order}.
\]

Это означает: один покупатель может иметь много заказов.

В Java это может быть представлено так:

\begin{verbatim}
public class Customer {
    private List<Order> orders;
}
\end{verbatim}

\subsection{Алгоритм обратного проектирования кода в UML-модель}

\textbf{Вход:} исходный код.

\textbf{Выход:} UML-модель.

\subsubsection*{Шаги}

\begin{enumerate}
    \item Выполнить лексический и синтаксический анализ исходного кода.
    \item Найти объявления классов, интерфейсов, перечислений.
    \item Для каждого класса создать UML-элемент \texttt{Class}.
    \item Для каждого интерфейса создать UML-элемент \texttt{Interface}.
    \item Преобразовать поля классов в UML-атрибуты.
    \item Преобразовать методы в UML-операции.
    \item Найти отношения наследования.
    \item Найти реализации интерфейсов.
    \item Найти ассоциации по полям ссылочных типов.
    \item Построить диаграмму классов.
    \item При необходимости применить автоматическую раскладку элементов.
\end{enumerate}

\subsubsection*{Пример}

Код:

\begin{verbatim}
public interface Repository<T> {
    T findById(int id);
}

public class UserRepository implements Repository<User> {
    public User findById(int id) {
        return null;
    }
}
\end{verbatim}

После обратного проектирования:
\begin{itemize}
    \item появляется интерфейс \texttt{Repository<T>};
    \item появляется класс \texttt{UserRepository};
    \item появляется отношение реализации:
    \[
    \texttt{UserRepository} \dashrightarrow \texttt{Repository<User>};
    \]
    \item операция \texttt{findById(id: int): User} отображается в модели.
\end{itemize}

\section{Примеры известных UML CASE-систем}

\subsection{Sparx Systems Enterprise Architect}

\textbf{Enterprise Architect} --- одна из наиболее известных промышленных UML CASE-систем.

Основные возможности:
\begin{itemize}
    \item поддержка UML;
    \item поддержка SysML, BPMN, ArchiMate и других нотаций;
    \item моделирование требований;
    \item построение диаграмм классов, последовательностей, состояний, активностей, компонентов и развёртывания;
    \item генерация документации;
    \item генерация и обратное проектирование кода;
    \item поддержка командной работы;
    \item работа с репозиториями моделей;
    \item трассировка требований;
    \item интеграция с внешними инструментами.
\end{itemize}

Enterprise Architect часто применяется в:
\begin{itemize}
    \item системной инженерии;
    \item разработке корпоративных информационных систем;
    \item архитектурном моделировании;
    \item моделировании бизнес-процессов;
    \item проектировании встроенных и технических систем.
\end{itemize}

\subsection{Visual Paradigm}

\textbf{Visual Paradigm} --- UML CASE-система и платформа визуального моделирования.

Основные возможности:
\begin{itemize}
    \item построение UML-диаграмм;
    \item поддержка диаграмм вариантов использования, классов, последовательностей, активностей, компонентов и развёртывания;
    \item моделирование бизнес-процессов;
    \item генерация документации;
    \item поддержка командной работы;
    \item инженерия кода;
    \item проектирование баз данных;
    \item поддержка требований и пользовательских историй.
\end{itemize}

Visual Paradigm удобен для учебных и промышленных проектов, где требуется не только рисование диаграмм, но и связанная модель, документация и командная работа.

\subsection{IBM Rational Software Architect}

\textbf{IBM Rational Software Architect} --- промышленная среда моделирования и проектирования, связанная с линейкой IBM Rational.

Основные возможности:
\begin{itemize}
    \item UML-моделирование;
    \item проектирование архитектуры приложений;
    \item интеграция с процессами корпоративной разработки;
    \item анализ моделей;
    \item проверка соответствия стандартам;
    \item поддержка шаблонов проектирования;
    \item интеграция с другими инструментами IBM.
\end{itemize}

Исторически семейство IBM Rational сыграло большую роль в распространении UML и объектно-ориентированного анализа и проектирования.

\subsection{IBM Rational Rose}

\textbf{Rational Rose} --- одна из классических CASE-систем для объектно-ориентированного анализа и проектирования.

Особенности:
\begin{itemize}
    \item поддержка UML-диаграмм;
    \item проектирование классов и взаимодействий;
    \item генерация кода;
    \item обратное проектирование;
    \item использование в методологиях объектно-ориентированного проектирования.
\end{itemize}

Rational Rose считается исторически важной системой: она была широко распространена в период активного становления UML как промышленного стандарта.

\subsection{Modelio}

\textbf{Modelio} --- UML/BPMN CASE-система, имеющая open-source-версию.

Основные возможности:
\begin{itemize}
    \item UML-моделирование;
    \item BPMN-моделирование;
    \item поддержка расширений и модулей;
    \item генерация документации;
    \item поддержка архитектурного моделирования;
    \item возможность расширения под конкретные процессы разработки.
\end{itemize}

Modelio часто рассматривается как доступная альтернатива коммерческим CASE-системам.

\subsection{StarUML}

\textbf{StarUML} --- кроссплатформенный инструмент UML-моделирования.

Основные возможности:
\begin{itemize}
    \item поддержка UML-диаграмм;
    \item поддержка расширений;
    \item генерация кода;
    \item обратное проектирование;
    \item поддержка SysML и ERD в современных версиях;
    \item удобство для небольших и средних проектов.
\end{itemize}

StarUML часто используется в учебных целях и в проектах, где требуется сравнительно лёгкий инструмент для проектирования.

\subsection{PlantUML}

\textbf{PlantUML} --- текстовый инструмент описания UML-диаграмм.

PlantUML не является классической CASE-системой в полном смысле, так как его основная функция --- генерация диаграмм из текстового описания. Однако он широко используется как средство UML-документирования.

Пример описания диаграммы классов на PlantUML:

\begin{verbatim}
@startuml
class User {
  - id: int
  - name: String
  + getName(): String
}

class Order {
  - number: String
}

User "1" -- "0..*" Order
@enduml
\end{verbatim}

Преимущества PlantUML:
\begin{itemize}
    \item диаграммы хранятся как текст;
    \item удобно использовать с Git;
    \item легко включать в документацию;
    \item хорошо подходит для continuous documentation;
    \item можно автоматически генерировать диаграммы при сборке проекта.
\end{itemize}

Ограничение PlantUML состоит в том, что он больше ориентирован на описание диаграмм, чем на ведение полноценного репозитория модели.

\section{Сравнение UML CASE-систем}

\begin{longtable}{p{0.18\textwidth}p{0.18\textwidth}p{0.18\textwidth}p{0.18\textwidth}p{0.18\textwidth}}
\toprule
\textbf{Система} & \textbf{UML-моделирование} & \textbf{Генерация кода} & \textbf{Командная работа} & \textbf{Типичное применение} \\
\midrule
Enterprise Architect &
Да &
Да &
Да &
Промышленные проекты, системная инженерия, архитектура \\
\addlinespace

Visual Paradigm &
Да &
Да &
Да &
Проектирование ПО, обучение, бизнес-моделирование \\
\addlinespace

IBM Rational Software Architect &
Да &
Да &
Да &
Корпоративная разработка, архитектура приложений \\
\addlinespace

Modelio &
Да &
Да, через модули &
Ограниченно/зависит от версии &
UML/BPMN-моделирование, open-source-проекты \\
\addlinespace

StarUML &
Да &
Да &
Ограниченно/через расширения &
Учебные и средние проекты \\
\addlinespace

PlantUML &
Частично, через текстовые диаграммы &
Нет как основная функция &
Через Git &
Документация, lightweight-моделирование \\
\bottomrule
\end{longtable}

\section{Преимущества использования UML CASE-систем}

\subsection{Повышение наглядности}

UML-диаграммы позволяют представить сложную систему в виде набора визуальных моделей. Это облегчает коммуникацию между:
\begin{itemize}
    \item аналитиками;
    \item архитекторами;
    \item разработчиками;
    \item тестировщиками;
    \item заказчиками;
    \item менеджерами проекта.
\end{itemize}

\subsection{Снижение числа ошибок на ранних стадиях}

Ошибки, обнаруженные на стадии проектирования, обычно дешевле исправлять, чем ошибки, обнаруженные после реализации. CASE-система помогает обнаруживать:
\begin{itemize}
    \item неполные требования;
    \item противоречивые сценарии;
    \item циклические зависимости;
    \item несогласованность интерфейсов;
    \item избыточные классы;
    \item отсутствие связей между требованиями и реализацией.
\end{itemize}

\subsection{Автоматизация рутинных операций}

CASE-системы автоматизируют:
\begin{itemize}
    \item создание шаблонов кода;
    \item построение документации;
    \item проверку модели;
    \item построение отчётов;
    \item обновление диаграмм;
    \item анализ зависимостей.
\end{itemize}

\subsection{Поддержка единого архитектурного описания}

В крупных системах важно иметь единый источник архитектурной информации. UML CASE-система может играть роль такого источника, если модель поддерживается в актуальном состоянии.

\section{Недостатки и ограничения UML CASE-систем}

\subsection{Сложность внедрения}

Полноценная CASE-система требует:
\begin{itemize}
    \item обучения сотрудников;
    \item настройки процесса;
    \item выбора уровня детализации моделей;
    \item интеграции с существующими инструментами;
    \item дисциплины ведения модели.
\end{itemize}

\subsection{Риск избыточного моделирования}

Если моделировать слишком подробно, модель может стать дороже самой разработки. Особенно это характерно для проектов с быстро меняющимися требованиями.

Рациональный подход:
\begin{itemize}
    \item моделировать архитектурно значимые части;
    \item не дублировать в UML очевидный код;
    \item использовать диаграммы для коммуникации и анализа;
    \item поддерживать только полезные модели.
\end{itemize}

\subsection{Проблема устаревания моделей}

Если модель не синхронизируется с кодом, она быстро теряет ценность. Поэтому важны:
\begin{itemize}
    \item round-trip engineering;
    \item автоматическое построение моделей по коду;
    \item интеграция с CI/CD;
    \item регулярный архитектурный контроль.
\end{itemize}

\subsection{Ограниченность генерации кода}

Генерация кода хорошо работает для:
\begin{itemize}
    \item каркасов классов;
    \item интерфейсов;
    \item DTO;
    \item схем данных;
    \item типовых шаблонов.
\end{itemize}

Но бизнес-логика обычно требует ручной реализации. Поэтому UML CASE-системы не заменяют программиста, а помогают структурировать разработку.

\section{UML CASE-системы и модельно-ориентированная разработка}

\subsection{MDA}

\textbf{MDA} --- \textit{Model Driven Architecture}, архитектура, управляемая моделями.

Идея MDA состоит в том, что основным артефактом разработки является модель, а код и другие реализации могут автоматически порождаться из неё.

В MDA различают:
\begin{itemize}
    \item \textbf{CIM} --- Computation Independent Model, модель, независимая от вычислительной реализации;
    \item \textbf{PIM} --- Platform Independent Model, модель, независимая от платформы;
    \item \textbf{PSM} --- Platform Specific Model, модель, специфичная для платформы;
    \item \textbf{Code} --- исходный код.
\end{itemize}

Типовая цепочка:
\[
\text{CIM} \rightarrow \text{PIM} \rightarrow \text{PSM} \rightarrow \text{Code}.
\]

\subsection{Пример}

Пусть есть независимая от платформы сущность:

\[
\texttt{Customer(id: Integer, name: String)}.
\]

Для Java/JPA она может быть преобразована в:

\begin{verbatim}
@Entity
public class Customer {
    @Id
    private Integer id;

    private String name;
}
\end{verbatim}

Для SQL она может быть преобразована в:

\begin{verbatim}
CREATE TABLE customer (
    id INTEGER PRIMARY KEY,
    name VARCHAR(255)
);
\end{verbatim}

Таким образом, одна модель может служить источником для нескольких реализаций.

\section{Практические рекомендации по применению UML CASE-систем}

\subsection{Когда UML CASE-система особенно полезна}

UML CASE-системы оправданы, если:
\begin{itemize}
    \item система большая и долгоживущая;
    \item требуется строгая документация;
    \item проект разрабатывается большой командой;
    \item важна трассировка требований;
    \item требуется архитектурное согласование;
    \item система имеет сложные состояния или взаимодействия;
    \item нужно генерировать документацию или код;
    \item система относится к критическим областям.
\end{itemize}

\subsection{Когда достаточно лёгких средств}

В небольших проектах часто достаточно:
\begin{itemize}
    \item PlantUML;
    \item Mermaid;
    \item draw.io;
    \item диаграмм в документации;
    \item архитектурных схем C4.
\end{itemize}

Если проект небольшой, а команда хорошо понимает архитектуру, тяжёлая CASE-система может оказаться избыточной.

\subsection{Оптимальный уровень детализации}

Не следует моделировать всё подряд. Лучше выбирать диаграммы по назначению:

\begin{longtable}{p{0.35\textwidth}p{0.55\textwidth}}
\toprule
\textbf{Задача} & \textbf{Подходящая диаграмма} \\
\midrule
Описание требований & Диаграмма вариантов использования \\
Описание предметной области & Диаграмма классов \\
Описание алгоритма или процесса & Диаграмма активностей \\
Описание взаимодействия объектов & Диаграмма последовательностей \\
Описание жизненного цикла объекта & Диаграмма состояний \\
Описание физической архитектуры & Диаграмма развёртывания \\
Описание модулей системы & Диаграмма компонентов \\
\bottomrule
\end{longtable}

\section{Краткий итог}

UML CASE-системы --- это инструменты автоматизированного анализа, проектирования, документирования и сопровождения программных систем на основе UML-моделей.

Их ключевые функции:
\begin{itemize}
    \item создание UML-диаграмм;
    \item хранение единой модели;
    \item проверка согласованности;
    \item трассировка требований;
    \item генерация документации;
    \item генерация кода;
    \item обратное проектирование;
    \item поддержка командной работы;
    \item интеграция с IDE и системами контроля версий;
    \item поддержка профилей и расширений UML.
\end{itemize}

Наиболее известные UML CASE-системы:
\begin{itemize}
    \item Sparx Systems Enterprise Architect;
    \item Visual Paradigm;
    \item IBM Rational Software Architect;
    \item IBM Rational Rose;
    \item Modelio;
    \item StarUML;
    \item PlantUML как лёгкое текстовое средство UML-документирования.
\end{itemize}

Главная ценность UML CASE-систем состоит не в красивом рисовании диаграмм, а в поддержке связанной, проверяемой и сопровождаемой модели системы.

\end{document}
